\documentclass[11pt,a4paper]{article}
\usepackage[T1]{fontenc}
\usepackage[utf8]{inputenc}
\usepackage{amsmath,amssymb,amsthm}
\usepackage[margin=1in]{geometry}
\usepackage[colorlinks=true,linkcolor=blue,urlcolor=blue]{hyperref}
\theoremstyle{plain}
\newtheorem{theorem}{Theorem}
\newtheorem{lemma}[theorem]{Lemma}
\newtheorem{proposition}[theorem]{Proposition}
\newtheorem{corollary}[theorem]{Corollary}
\theoremstyle{definition}
\newtheorem{remark}[theorem]{Remark}

\title{The empirical recurrence for OEIS A224743, proved by transfer matrix}
\author{Adrian Perez Fontelles\\ \small Independent researcher}
\date{2 September 2026}
\begin{document}
\maketitle

\begin{abstract}
OEIS A224743 counts the $(n+1)\times7$ matrices over $\{0,\dots,1\}$ all of whose
$2\times2$ contiguous subblocks are of equal permanent, and carries an empirical recurrence of order
$34$ contributed by R.~H.~Hardin. It is true, and it is decidable rather than empirical.
The condition ties together $2$ consecutive rows and $2$ consecutive columns at once, so a
single row is not a state; the strip of the last $1$ rows is, and the admissible strips
are read off the overlap graph of the admissible windows instead of being enumerated. That
makes $a(n)$ a walk count on $S=163$ vertices, hence $C$-finite, and whether the conjectured
recurrence annihilates it is settled by a finite exact computation.
\end{abstract}

\noindent\small 2020 Mathematics Subject Classification. 05A15, 05B45, 15B36.\normalsize

\section{The conjecture}

OEIS A224743 is ``Number of (n+1)X7 0..1 matrices with each 2X2 permanent equal.''. It has offset $1$ and begins
\[
1349, 40297, 661569, 15845297, 316084673, 6987578401, 146736389761, 3168731910657,\ \dots
\]
The entry states, as a conjecture contributed by its author and never marked settled:
\begin{quote}
Empirical: a(n) = 34*a(n-1) -75*a(n-2) -7140*a(n-3) +54718*a(n-4) +428776*a(n-5) -5148288*a(n-6) -7894384*a(n-7) +216431848*a(n-8) -145979616*a(n-9) -5025133232*a(n-10) +9340472128*a(n-11) +69509336640*a(n-12) -190959984640*a(n-13) -579126695808*a(n-14) +2166687694336*a(n-15) +2696767474176*a(n-16) -15137991368704*a(n-17) -4263323315200*a(n-18) +66643415601152*a(n-19) -21968505323520*a(n-20) -181114010091520*a(n-21) +142147429171200*a(n-22) +282120150515712*a(n-23) -347524113432576*a(n-24) -202500951769088*a(n-25) +419006127800320*a(n-26) -2311631929344*a(n-27) -237796008656896*a(n-28) +74252761956352*a(n-29) +49986725216256*a(n-30) -26211312992256*a(n-31) -1579608440832*a(n-32) +2358473916416*a(n-33) -274877906944*a(n-34)
\end{quote}
As of the ``Last modified'' line on the live entry (Oct 03 2025 23:36:19, revision 8) this is still
recorded as empirical, and nothing on the entry records it as proved.

\section{The count is a walk count}

Call a $2\times2$ matrix over $\{0,\dots,1\}$ a \emph{window} when it is of equal permanent, and
call a $2\times7$ matrix a \emph{band} when every one of its $6$ contiguous
$2\times2$ subblocks is a window. What the entry counts is the matrices of $7$ columns
whose every $2$ consecutive rows form a band.

Take as states the $1\times7$ strips that occur as the top $1$ rows of a band,
together with those that occur as its bottom $1$ rows, and put an edge from $u$ to $v$
when the band whose first $1$ rows are $u$ and whose last $1$ rows are $v$ exists.
A band is determined by that pair, so the edge is unique when it exists. Reading a counted
matrix from the top, its rows $1..1$, then $2..2$, and so on, form a walk, and every
walk arises from exactly one matrix; a matrix with $n+1$ rows gives a walk of $n$ edges.
Since the first $1$ rows of a counted matrix with at least $2$ rows are the top of a
band, and since no further condition falls on where a walk starts or ends,
\[
a(n)\;=\;\iota^{\!\top}M^{\,n}\tau ,\qquad \iota=\tau=(1,1,\dots,1)^{\!\top},
\]
with $M$ the adjacency matrix on $S=163$ states. In particular $a$ is $C$-finite. The
alignment was checked against the entry's published terms rather than assumed.

There are $(1 + 1)^{7}$ strips of $1$ rows and $7$ columns, far too
many to test one by one, and they are not tested one by one. Two consecutive windows of a band
overlap in $2\times1$ entries, so a band is exactly a walk of $6$ windows in
the graph that joins a window to those whose first $1$ columns are its last $1$
columns. Enumerating those walks produces the bands, and with them the states and the edges,
directly.

\section{The admissible windows}

The entry does not name the common value; it only asks that all the windows of one matrix agree.
So fix a value $v$ and let $W_v$ be the set of $2\times2$ matrices over
$\{0,\dots,1\}$ whose permanent, $b_{11}b_{22}+b_{12}b_{21}$, equals $v$. A matrix counted by the entry has all of
its windows in a single $W_v$, and the value of $v$ is determined by the matrix, so the sets of
matrices belonging to different $v$ are disjoint and
\[
a(n)\;=\;\sum_v a_v(n),
\]
each $a_v$ being the walk count of Section~2 built from $W_v$ alone. A disjoint union of the
$3$ digraphs, one per value of $v$, therefore counts the whole of $a(n)$ with the same
all-ones vectors, and the $2\times2$ windows in play number $16$ in total.

\section{The proof}

\begin{lemma}\label{lem:crit}
Let $q(t)=t^{34}-\sum_i c_it^{34-i}$ be the characteristic polynomial of the
conjectured recurrence. The residual $a(n)-\sum_ic_ia(n-i)$ equals
$\iota^{\!\top}M^{\,j}q(M)\tau$ for the corresponding walk index $j$, so the recurrence
holds from some point on if and only if $\iota^{\!\top}M^{j}q(M)\tau=0$ for all large $j$,
and it is enough to examine $j<S$.
\end{lemma}

\begin{proof}
Factor $M^{j}$ out of $M^{j+34}-\sum_ic_iM^{j+34-i}$. For the bound, $M^{S}$
is an integer combination of $I,M,\dots,M^{S-1}$ by Cayley--Hamilton, so the numbers
$u_j=\iota^{\!\top}M^{j}q(M)\tau$ obey the monic recurrence given by the characteristic
polynomial of $M$; $S$ consecutive zeros force every later one, and the last nonzero $u_j$
pins the threshold exactly.
\end{proof}

\begin{theorem}
\[
a(n)\;=\;34\,a(n-1) - 75\,a(n-2) - 7140\,a(n-3) + 54718\,a(n-4) + 428776\,a(n-5) - 5148288\,a(n-6) - 7894384\,a(n-7) + 216431848\,a(n-8) - 145979616\,a(n-9) - 5025133232\,a(n-10) + 9340472128\,a(n-11) + 69509336640\,a(n-12) - 190959984640\,a(n-13) - 579126695808\,a(n-14) + 2166687694336\,a(n-15) + 2696767474176\,a(n-16) - 15137991368704\,a(n-17) - 4263323315200\,a(n-18) + 66643415601152\,a(n-19) - 21968505323520\,a(n-20) - 181114010091520\,a(n-21) + 142147429171200\,a(n-22) + 282120150515712\,a(n-23) - 347524113432576\,a(n-24) - 202500951769088\,a(n-25) + 419006127800320\,a(n-26) - 2311631929344\,a(n-27) - 237796008656896\,a(n-28) + 74252761956352\,a(n-29) + 49986725216256\,a(n-30) - 26211312992256\,a(n-31) - 1579608440832\,a(n-32) + 2358473916416\,a(n-33) - 274877906944\,a(n-34)
\]
for every $n>34$.
\end{theorem}

\begin{proof}
The vector $w=q(M)\tau$ was formed by $34$ matrix--vector products in exact integer
arithmetic and $\iota^{\!\top}M^{j}w$ evaluated for $j=0,1,\dots$, again exactly; those
integers vanish from the index corresponding to $n=34+1$ onwards and the run continued
until the working vector was identically zero, which settles every larger $j$ at once.
\end{proof}

\section{Verification}

Three checks, each independent of the algebra above.

First, the model reproduces all $14$ terms the entry publishes, exactly, in integer
arithmetic. This is what ties the model to the entry: the digraph is built by reading the
entry's English, and a misreading gives different counts.

Second, the conjectured recurrence was evaluated directly on those published terms, with no
matrices involved, and holds wherever the proved range and the data overlap.

Third, the annihilation test of Lemma~\ref{lem:crit} was carried out in exact integer
arithmetic on the whole vector rather than on a sampled prefix.

\begin{thebibliography}{9}
\bibitem{oeis} The OEIS Foundation, \emph{The On-Line Encyclopedia of Integer
Sequences}, \url{https://oeis.org}, 2026. Sequence A224743.
\bibitem{stanley} R.~P.~Stanley, \emph{Enumerative Combinatorics}, Volume 1, 2nd ed.,
Cambridge University Press, 2011. (Transfer-matrix method, Section 4.7.)
\bibitem{kim} K.~H.~Kim, \emph{Boolean Matrix Theory and Applications}, Marcel Dekker,
1982.
\bibitem{horn} R.~A.~Horn and C.~R.~Johnson, \emph{Matrix Analysis}, 2nd ed., Cambridge
University Press, 2013.
\end{thebibliography}

\end{document}
